\scanpage{609}\mspage{18}
% archive sheet number: 608

Donc il reste à prouver que tt $x\in\mathrm{Ob}\,\underline S^{\unsure{*|\prime}}$
est dans l'image essentielle. Or l'hypothèse nous
\deleted{\unsure{dit}}\inserted{assure} que
$x\simeq\varinjlim_{i\in\deleted{\unsure{\Phi}}I_x}x_i$ peut s'écrire
$\unsure{\varinjlim|\underline{\lim}}_{i\in I_{x,\underline S'}}x_i$
(car $I_{x,\underline S'}$ est cofinal dans $I_x$), où cette fois les
$x_i$ sont dans $\underline S'$. Comme $I_x$ \unsure{gros} devant
$\pi$, de $=$ la sous-catégorie pleine cofinale $I_{x\underline S}$,
donc on \unsure{gagne|y gagne}.

\uline{Corollaire} \deleted{\illegible{1}} Sous les conditions
précédentes, si $\underline S'$ est $\pi$-accessible, il en est de
$=$ de l'ouvert complémentaire $\underline S$, \uline{pourvu} que
$\underline S$ soit stable \inserted{dans $M$} par $\varinjlim$ de type
$L_\pi$\margin{!}

Il faut prouver que $\underline S$ est égale à l'image essentielle
de $\mathrm{Ind}_{\pi}\,\underline S(\pi)$.
\deleted{Soit donc}\inserted{L'hypothèse ``\unsure{pourvu}''}
\deleted{Prouvons}\inserted{assure} \deleted{d'abord}
que cette image essentielle est $\subset\underline S$.
\deleted{\unsure{Inversement}, soit donc
$\left|\,x=\varinjlim_I x_i\quad
\left|\begin{aligned}&x_i\in\mathrm{Ob}\,\underline S(\pi)\\
&I\ \text{\unsure{gros} devant }\pi\end{aligned}\right.\right)$,}
prouvons que \deleted{prouvons que $x\in\mathrm{Ob}\,\underline S$.
Sinon, on aurait} \deleted{\unsure{pour}}
$x\in\mathrm{Ob}\,\underline S$ est dans cette image essentielle.
On aura
\[
 x\simeq\varinjlim_{i\in I}x_i\qquad
 \left|\begin{aligned}&x_i\in\mathrm{Ob}\,M_\pi\\
 &I=M_\pi/x\\&\text{\unsure{gros} devant }\pi\end{aligned}\right.,
\]
et comme $x\in\mathrm{Ob}\,\underline S$ et $\underline S$ ouvert,
les $x_i$ \deleted{prouvons que $\varinjlim x_i\in$} sont dans
$\underline S$, donc dans $\underline S(\pi)$, OK.

\deleted{\uline{\unsure{Remarque}}} Dans le cas qui nous occupe,
\inserted{où $M=(k\text{-alg})$,} comme $\underline S$ n'est
justement pas stable dans $M$ pour $L_\pi$
\inserted{(quel que soit $\pi$)}, \unsure{(c'est|c'est} \unsure{cela qui|cela que} nous faisait conclure que
